\chapter{Correcciones y confirmacion manual}
\label{cap:correcciones-manuales}

\begin{procedimiento}{Objetivo del capitulo}{Operador del torneo}
Explicar como corregir resultados ya guardados, cuando cancelar una confirmacion sensible y como usar el modal de participante solo como herramienta de soporte controlado.
\end{procedimiento}

Este capitulo no reemplaza los procedimientos de grupos, batallas ni fases manuales. Se usa cuando ya existe un resultado persistido y el operador necesita corregirlo despues de revisar el impacto sobre fases posteriores.

\section{Cuando una correccion es sensible}

\begin{tabularx}{\linewidth}{>{\bfseries}p{3.8cm}X}
\toprule
Situacion & Criterio operativo\\
\midrule
Seleccion sin guardar & Puede ajustarse desde la misma pantalla antes de presionar el boton de guardado. No hay cambio persistido todavia.\\
Resultado guardado en la fase actual & Debe corregirse desde la misma tarjeta de grupo o batalla y volver a guardar. Revise que el mensaje de exito aparezca al final.\\
Resultado guardado con fases posteriores iniciadas & Es una operacion sensible. La plataforma solicita una confirmacion adicional porque la correccion puede afectar cruces, estados o historial ya generados.\\
Participante con estado incoherente & Revise el modal de participante antes de editar. Use la edicion manual solo si la organizacion identifica una inconsistencia concreta.\\
\bottomrule
\end{tabularx}

\begin{operacionsensible}
No aplique correcciones durante una prueba, demostracion publica o transmision del torneo sin autorizacion de la mesa. Antes de aceptar una correccion, anote la fase, el grupo o batalla, el resultado anterior, el resultado nuevo y la razon operativa.
\end{operacionsensible}

\section{Corregir un resultado guardado}

\figuraManual{capturas/finales/MU-028-confirmacion-correccion.png}{Estado previo a una confirmacion de correccion sensible en una fase posterior.}{fig:bloque4-mu-028}

\begin{procedimiento}{Aplicar una correccion de resultado}{Operador del torneo}
Use este procedimiento para cambiar un clasificado de grupo, un ganador de batalla o los clasificados multiples de una batalla cuando el resultado anterior ya fue guardado.
\end{procedimiento}

\begin{precondiciones}
\begin{itemize}
    \item El operador esta autenticado en el panel interno.
    \item El resultado oficial corregido fue validado por la organizacion.
    \item Se conoce que fase, grupo o batalla contiene el dato equivocado.
    \item La pantalla muestra los equipos ficticios o reales que corresponden a la operacion actual.
\end{itemize}
\end{precondiciones}

\paso{1}{Abra la division y la fase donde esta el resultado a corregir.}
\paso{2}{Ubique el grupo o la batalla afectada.}
\paso{3}{Compare la seleccion guardada con el acta o decision oficial.}
\paso{4}{Cambie la seleccion en la tarjeta correspondiente: \botoninterfaz{Marcar ganador}, \botoninterfaz{Clasificar} o el control visible equivalente.}
\paso{5}{Revise el contador o el ganador resaltado antes de guardar.}
\paso{6}{Presione \botoninterfaz{Guardar este grupo}, \botoninterfaz{Guardar este ganador}, \botoninterfaz{Guardar clasificados} o \botoninterfaz{Guardar todos los resultados}, segun el caso.}
\paso{7}{Si aparece la confirmacion de guardado, aceptela solo despues de revisar la seleccion visible.}
\paso{8}{Si aparece la confirmacion adicional de correccion, lea el texto completo y decida aceptar o cancelar.}
\paso{9}{Espere el mensaje de exito y vuelva a revisar la fase posterior afectada.}
\paso{10}{Registre la correccion en la bitacora operativa del evento.}

\begin{operacionsensible}
\textbf{Confirmación nativa de guardado de resultados.}
La apariencia exacta del diálogo depende del navegador y del sistema operativo. Texto exacto documentado:
\emph{¿Seguro que deseas guardar estas selecciones? Podrás editarlas después, pero verifica bien antes de continuar.}

Opciones reales: aceptar o cancelar. Al aceptar, la operación se envía al servidor. Al cancelar, no se guardan los cambios. Recomendación: revisar selección, grupo y cantidad de clasificados antes de aceptar.
\end{operacionsensible}


\begin{operacionsensible}
\textbf{Confirmación nativa de corrección manual.}
La apariencia exacta del diálogo depende del navegador y del sistema operativo. Texto exacto documentado:
\emph{Este cambio puede afectar fases posteriores ya iniciadas. ¿Deseas aplicar la corrección manual?}

Opciones reales: aceptar o cancelar. Al aceptar, se envía la corrección con \texttt{manual\_override}. Al cancelar, se conserva el resultado anterior. Recomendación: usarlo solo como contingencia operativa documentada.
\end{operacionsensible}


\begin{resultadoesperado}
La plataforma muestra \campointerfaz{Correccion aplicada. Revisa las fases posteriores afectadas.} cuando la correccion sensible fue aceptada y procesada por el servidor.
\end{resultadoesperado}

\begin{fallas}
Si cancela el dialogo, el resultado anterior se conserva. Si no aparece mensaje de exito, no avance: recargue la fase, revise cual resultado quedo guardado y repita el procedimiento solo con la informacion confirmada.
\end{fallas}

\section{Correcciones por tipo de resultado}

\begin{tabularx}{\linewidth}{>{\bfseries}p{3.6cm}X}
\toprule
Tipo & Revision antes de guardar\\
\midrule
Clasificados de grupo & El contador debe coincidir con la cantidad requerida para ese grupo. Si el grupo pide dos clasificados, no deje uno ni tres.\\
Ganador unico & Solo un equipo debe quedar marcado como ganador. Revise que el perdedor visible corresponda al resultado oficial.\\
Clasificados multiples & La etiqueta \campointerfaz{Clasifican n} define la cantidad exacta. No use la memoria de otra fase como referencia.\\
Guardado general & Revise todos los formularios modificados visibles, porque el boton global procesa mas de una tarjeta si hay cambios pendientes.\\
\bottomrule
\end{tabularx}

\section{Modal de participante}

\figuraManual{capturas/finales/MU-029-modal-participante.png}{Modal de participante con datos de estado, ubicacion actual e historial.}{fig:bloque4-mu-029}

El modal de participante ayuda a verificar donde esta un equipo dentro de la competencia. Sirve para consultar fase, grupo, batalla, estado actual e historial antes de decidir si una correccion es necesaria.

\begin{procedimiento}{Consultar el detalle de un participante}{Operador del torneo}
Use este recorrido antes de editar manualmente el estado de un equipo.
\end{procedimiento}

\paso{1}{Abra la fase donde aparece el participante.}
\paso{2}{Pulse la tarjeta o control de detalle del participante si esta disponible.}
\paso{3}{Revise el resumen del equipo, su fase actual, grupo o batalla actual y estado.}
\paso{4}{Lea el historial para identificar el ultimo cambio registrado.}
\paso{5}{Cierre el modal si la informacion coincide con la competencia real.}

\begin{resultadoesperado}
El operador comprende el estado actual del participante sin modificar datos.
\end{resultadoesperado}

\section{Edicion manual de participante}

\figuraManual{capturas/finales/MU-030-edicion-participante.png}{Campos editables del modal de participante para soporte operativo controlado.}{fig:bloque4-mu-030}

\begin{procedimiento}{Editar estado de participante}{Operador del torneo}
Use esta accion solo para una correccion puntual aprobada. No la use para operar resultados normales.
\end{procedimiento}

\paso{1}{Abra el modal del participante y confirme que se trata del equipo correcto.}
\paso{2}{Revise la fase actual, el grupo actual, la batalla actual y el estado manual.}
\paso{3}{Cambie solo el campo necesario para resolver la inconsistencia.}
\paso{4}{Escriba una nota breve con la razon de la correccion si el formulario la solicita.}
\paso{5}{Presione \botoninterfaz{Guardar cambios}.}
\paso{6}{Verifique que el modal y la fase muestran el estado actualizado.}
\paso{7}{Revise cualquier fase posterior que dependiera de ese participante.}

\begin{advertencia}
Mover un participante a una fase, grupo o batalla incompatible puede dejar la competencia dificil de interpretar para otros operadores. Si el destino no es claro, cancele y revise primero desde Django Admin o con el responsable tecnico.
\end{advertencia}

\section{Tabla de recuperacion}

\begin{tabularx}{\linewidth}{>{\bfseries}p{4.2cm}X}
\toprule
Problema & Accion recomendada\\
\midrule
Se selecciono el equipo equivocado antes de guardar & Corrija la seleccion en pantalla y guarde solo cuando el contador o ganador sea correcto.\\
Se acepto una correccion equivocada & No haga cambios encadenados. Documente el error, recargue la fase y determine el ultimo estado persistido antes de corregir de nuevo.\\
No aparece la confirmacion sensible & Puede que aun no existan fases posteriores afectadas. Aun asi, revise el resultado y el mensaje de exito.\\
El servidor devuelve error de validacion & Ajuste la cantidad de clasificados, ganador o destino indicado por el mensaje.\\
La fase posterior quedo desactualizada & Revise sus batallas, participantes vivos y eliminados antes de operar nuevos resultados.\\
Hay duda sobre el impacto & Cancele la operacion. Una correccion sensible cancelada conserva el estado anterior.\\
\bottomrule
\end{tabularx}

\section{Checklist del capitulo}

\begin{itemize}
    \item Resultado anterior identificado por fase, grupo o batalla.
    \item Resultado nuevo validado por la organizacion.
    \item Cantidad de clasificados o ganador unico revisado antes de guardar.
    \item Confirmacion de guardado leida antes de aceptar.
    \item Confirmacion adicional aceptada solo si el impacto fue entendido.
    \item Fases posteriores revisadas despues de la correccion.
    \item Correccion registrada fuera de la plataforma si el evento lo exige.
\end{itemize}
